% 최종 축약본. chapter4.md와 같은 본문에서 생성.
\subsection{재정렬 모델 교체에 따른 근거 선택 성능}
\label{sec:ch4_evidence_results}

Jev의 근거 F1은 61.75점으로 CE보다 3.19점 높았고, Recall은 2.95점 높았다(표~\ref{tab:ch4_main_results}). F1은 정답 근거의 누락과 정답 외 근거의 선택을 함께 반영하며, Recall은 정답 근거 중 회수한 비율이다. 두 점수는 문항별로 계산한 뒤 같은 가중치로 평균하였다. 정답 근거는 평가 자료가 지정한 법령이며, 선택한 근거와 문자열이 같은지를 비교하였다. 따라서 정답과 다른 근거가 법률적으로 타당한지까지 판정하는 점수는 아니다.

\begin{table}[htbp]
\centering
\caption{전체 226문항의 성능. $\Delta$는 Jev$-$CE이다.}
\label{tab:ch4_main_results}
\small
\setlength{\tabcolsep}{3pt}
\begin{tabular}{@{}lrrrr@{}}
\hline
지표 & CE & Jev & $\Delta$ & 차이의 95\% CI \\
\hline
\multicolumn{5}{@{}l}{\textbf{근거 선택}} \\
F1 & 58.57 & 61.75 & +3.19 & [+0.75, +5.68] \\
Recall & 57.74 & 60.69 & +2.95 & [+0.39, +5.51] \\
EM & 27.88 & 29.65 & +1.77 & [-1.77, +5.19] \\
완전 회수율 & 34.07 & 34.96 & +0.88 & [-2.90, +4.53] \\
\multicolumn{5}{@{}l}{\textbf{답변 품질}} \\
Token F1 & 46.20 & 47.31 & +1.12 & [+0.02, +2.27] \\
LF-Eval & 67.57 & 70.15 & +2.58 & [+0.37, +4.79] \\
\hline
\end{tabular}
\par\smallskip
\begin{minipage}{\linewidth}\footnotesize
점수는 원값의 100배이다. EM·완전 회수율은 백분율이며 차이는 \%p, 나머지 차이는 점 단위이다. 차이와 신뢰구간(CI)은 반올림 전 값으로 계산하였다. 출처나 정답 근거를 공유하는 문항을 묶은 151개 그룹의 대응 bootstrap 10,000회로 탐색적 구간을 산출했으며, 다중 비교는 보정하지 않았다.
\end{minipage}
\end{table}

반면 정답 근거를 모두 확보한 문항은 77개에서 79개로 늘어나는 데 그쳤다. 완전 회수율은 정답 근거를 빠짐없이 선택한 비율이고, EM은 정답 외 근거도 없이 정확히 같은 집합을 선택한 비율이다. 완전 회수율의 차이는 0.88\%p였으며 신뢰구간은 [-2.90, +4.53]\%p였다. EM의 차이도 1.77\%p로 신뢰구간에 0을 포함하였다. 따라서 평균적인 회수 수준은 높아졌지만, 필요한 근거를 모두 갖춘 질문이 뚜렷하게 늘었다고 보기는 어렵다.

재정렬 전 두 조건이 공유한 검색 후보의 평균 정답 근거 Recall은 77.21\%였다. 재정렬 후 상위 10개 후보에 남은 정답 근거의 평균 Recall은 CE 69.47\%, Jev 74.63\%였다. 그러나 후보 목록마다 하나만 선택하는 제약 때문에, 상위 후보에 정답이 있어도 모두 선택할 수 있는 것은 아니다. 이 제약 아래 완전 회수가 가능한 문항은 각각 104개와 120개였으나 실제로는 77개와 79개에 머물렀다. 즉, 상위 후보의 개선이 최종 선택으로 모두 이어지지는 않았다. 최종 성능은 재정렬 모델과 후속 언어 모델 선택을 함께 거친 결과로 해석해야 한다.

공개 코드 CE의 F1은 48.72점이었다. CE와의 9.85점 차이는 후보 목록 참조 오류를 고친 결과로, Jev 교체의 개선폭과 구별된다. 공개 코드는 현재 후보 목록에서 고른 번호를 첫 번째 목록에 적용하는 문제가 있었으며, 수정 조건은 현재 목록을 참조하도록 고쳤다. 이 오류가 원 논문의 실제 실행에도 있었는지는 확인되지 않았다.
